% ------------------------------------------------------------------ 2. Sistema
\seccionnueva
\section{Sistema de software}
\label{sec:sistema}

\begin{minipage}[c]{0.2\textwidth}
  \centering\includegraphics[width=2.6cm]{figuras/collaboard-logo.png}
\end{minipage}\hfill
\begin{minipage}[c]{0.76\textwidth}
  {\sffamily\bfseries\LARGE CollaBoard}\par\vspace{0.2em}
  {\color{cbMuted}Plataforma de trabajo colaborativo visual en tiempo real: tableros
  compartidos para dibujar, conversar y coordinar equipos, más un lienzo colectivo de píxeles.}
\end{minipage}

\subsection{Descripción del alcance entregado}

En el prototipo 1 CollaBoard era un conjunto de aplicaciones de una sola página (SPA) que el
navegador conectaba directamente a cada servicio. El prototipo 2 la convierte en un
\textbf{sistema distribuido}: dos presentaciones distintas (una web con renderizado en el
servidor y un cliente local de escritorio) acceden por un único \textbf{API Gateway} a seis
servicios de negocio escritos en cinco lenguajes, con un componente de procesamiento asíncrono
conectado por \textbf{RabbitMQ} y persistencia separada por dominio.

\begin{table}[H]
\centering
\small
\filas
\begin{tabularx}{\linewidth}{@{}P{3.6cm} L P{2.6cm}@{}}
  \encabezado \hd{Capacidad} & \hd{Qué puede hacer la persona usuaria} & \hd{Estado} \\
  Cuenta y sesión & Registrarse, iniciar y cerrar sesión; la web guarda la sesión en una cookie HttpOnly & \hecho \\
  Tableros & Crear, renombrar, borrar y salir; filtrar y buscar; ver el dashboard renderizado en el servidor & \hecho \\
  Colaboración & Invitar por link con rol (editor o lector), cambiar roles, expulsar miembros & \pone \\
  Dibujo simultáneo & Lápiz, figuras, notas, borrador, deshacer; ver cursores y trazos ajenos en vivo & \hecho \\
  Guardado e historial & Autosave, estado restaurado al recargar, versiones y restauración & \integracion \\
  Chat por tablero & Historial paginado, envío sin duplicados, recepción en vivo, reintento & \integracion \\
  Bandeja de notificaciones & Avisos de invitación aceptada y de mensajes; marcar como leída & \integracion \\
  Pixels & Lienzo colectivo de 200×200 con enfriamiento de 5 s & \pone \\
  Cliente de escritorio & Login, tableros, dibujo, chat y bandeja contra el mismo gateway & \porhacer \\
  \bottomrule
\end{tabularx}
\caption{Alcance funcional del prototipo 2. \textit{En integración}: interfaz terminada y probada contra
mocks del contrato; falta conectar el servicio definitivo.}
\label{tab:alcance}
\end{table}

\begin{pendiente}{Sebastian (compila) — todos}
Actualizar la columna \emph{Estado} el 1 de noviembre con el resultado real de las pruebas de
integración (sección~\ref{sec:trazabilidad}). Ninguna fila debe quedar en \emph{Pendiente} ni
\emph{En integración} el día de la entrega.
\end{pendiente}

\subsection{Requisitos del prototipo 2 y cómo se cumplen}

\begin{table}[H]
\centering
\small
\filas
\begin{tabularx}{\linewidth}{@{}P{4.4cm} L@{}}
  \encabezado \hd{Requisito} & \hd{Cumplimiento en CollaBoard} \\
  Arquitectura distribuida & Cada API, la web SSR, el \textit{worker}, el broker y cada base de datos corren en contenedores separados \\
  Dos presentaciones, una web con SSR & Web SSR (Next.js, TypeScript) y cliente local (Electron, TypeScript) \\
  Al menos seis componentes lógicos & Auth, Colaboración, Pizarra, Chat, Pixels y NotificationAPI, más MessageWorker \\
  Comunicación entre componentes & HTTP interno de permisos (Pizarra/Chat → Colaboración) y de eventos (→ NotificationAPI) \\
  Al menos cuatro componentes de datos (relacional y NoSQL) & UserDB, BoardDB, CanvasDB y NotificationDB en PostgreSQL; PixelDB (con la base de Chat) en MongoDB \\
  Procesos asíncronos & NotificationAPI publica en RabbitMQ; MessageWorker consume con confirmación manual \\
  Conectores HTTP & REST entre clientes, gateway y APIs; BFF de la web \\
  Cinco lenguajes de propósito general & TypeScript, Python, Go, Java y C\# \\
  Despliegue orientado a contenedores & Docker Compose para servidores; el cliente local se empaqueta aparte \\
  \bottomrule
\end{tabularx}
\caption{Requisitos del enunciado y su correspondencia en el sistema.}
\label{tab:requisitos}
\end{table}

\begin{nota}[Criterio de conteo]
No se cuentan HTML, CSS, SQL, JSON, YAML, Dockerfile ni la configuración de Nginx como lenguajes;
el JavaScript que produce TypeScript tampoco se cuenta aparte. RabbitMQ no se cuenta como base de datos.
\end{nota}
